% Explain what has been done by the R&D community and similar applications.
% Compare strengths, weaknesses and novelty of your solution.
\subsubsection*{Exact Methods}

One of the earliest exact approaches proposed for solving the OP is the Branch-and-Cut algorithm introduced by Fischetti et al.\cite{fischetti_solving_1998}, which combines linear programming relaxation, branching, and cutting planes to reduce the search space. Although effective on small and medium-sized instances, its computational cost increases rapidly with the number of vertices. This motivated the development of heuristic approaches, notably those of Tsiligiridis\cite{Tsiligiridis} and Golden and Levy\cite{golden}, to obtain near-optimal solutions more efficiently.

For the OPTW, Kantor and Rosenwein\cite{77719d57-9aea-3d04-aeff-58a2183cbee5} proposed a heuristic tree-search algorithm that incrementally constructs routes by adding temporally feasible vertices while continuously verifying time-window constraints. Although effective only on small instances, the method highlighted the strong dependency between vertices induced by temporal constraints. Righini and Salani later\cite{RIGHINI20091191} proposed an exact approach reformulating the OPTW as a Resource Constrained Elementary Shortest Path Problem (RCESPP), solved via dynamic programming with Decremental State Space Relaxation (DSSR). Despite significantly reducing the number of explored states through dominance and relaxation strategies, the method still suffers from state-space explosion on large instances and remains sensitive to the structure of the time windows.

\subsubsection*{Metaheuristic Methods}

Classical local search methods remained limited by their tendency to become trapped in local optima, motivating the development of more advanced metaheuristics capable of balancing intensification and diversification. Among the most influential approaches for the OPTW is the Variable Neighborhood Search (VNS) introduced by Labadie et al.\cite{LABADIE201215}, which alternates between several neighborhood structures while perturbing and reconstructing the current solution to escape local minima. Although the method demonstrated strong performance and hybridization potential, it remains sensitive to perturbation strength and parameter tuning.

Tabu Search, proposed by Archetti et al.\cite{Archetti}, incorporates adaptive memory structures to prevent revisiting previously explored solutions and reduce cycling behavior. The method achieves effective intensification and broad exploration, though its performance strongly depends on the configuration of tabu tenure and diversification strategies. Simulated Annealing (SA) has also been widely applied to the OPTW; inspired by thermodynamic processes, it probabilistically accepts deteriorating solutions to escape local optima, gradually shifting from exploration to exploitation as the temperature decreases. Its effectiveness, however, remains highly dependent on the cooling schedule and parameter configuration.

The Greedy Randomised Adaptive Search Procedure (GRASP), introduced by Feo and Resende\cite{feo_greedy_1995}, combines a randomised greedy construction phase with an iterative local search phase. At each iteration, a candidate solution is built by selecting elements from a Restricted Candidate List (RCL) that balances quality and randomness, and is then improved by local search; the best solution across all iterations is returned. GRASP has been applied to several routing problems due to its simplicity and its ability to explore diverse regions of the solution space without maintaining a population, making it a lightweight alternative to population-based metaheuristics.

\subsubsection*{Population-Based and Learning-Based Methods}

Karbowska-Chili\'{n}ska and Zabielski\cite{Karbowska} proposed a Genetic Algorithm (GA) that evolves a population of candidate routes through crossover and mutation operators. Although GAs provide strong exploration capabilities, they often struggle with feasibility preservation under time-window constraints, sensitivity to parameter tuning, and premature convergence. Ant Colony Optimization (ACO), proposed by Montemanni and Gambardella\cite{montemanni2023antcolonyteamorienteering}, incrementally constructs routes using pheromone-guided probabilistic decisions, demonstrating strong constructive and cooperative search capabilities while remaining sensitive to parameterization and premature convergence toward suboptimal solutions.

More recent approaches have explored Reinforcement Learning (RL) techniques, learning routing policies directly from interactions with the optimization environment. Methods such as Pointer Networks and attention-based models trained with policy gradient algorithms have been applied to variants of the orienteering problem, enabling end-to-end learning of construction heuristics without explicit problem-specific engineering. Although promising, such methods generally require substantial computational resources, large training datasets, and still lack the robustness and constraint-handling precision of established hybrid metaheuristics on tightly constrained instances such as the OPTW.

\subsubsection*{This Work}

In this report, we present a comparative study of several algorithms for solving the OPTW in the context of touristic landmark visit planning in Algiers. The study focuses on implementing heuristic, metaheuristic, and exact approaches while introducing enhancements intended to overcome some of their known limitations. Specifically, we implement a Greedy algorithm extended with a profit-to-time-commitment ratio criterion that accounts for both travel time and visit duration, avoiding the over-selection of long-duration landmarks that standard greedy heuristics exhibit. We implement a Simulated Annealing algorithm comparing two cooling schedules (exponential and linear) alongside two acceptance criteria (Boltzmann and Cauchy), with a reheating mechanism to escape temperature collapse. We implement a Genetic Algorithm based on Order Crossover (OX) and introduce a tailored feasibility-preserving genetic operator based on the MaxShift quantity, which enables crossover and insertion moves to be validated without full route re-simulation. We also implement the CPLEX exact MILP solver with an MTZ-based time-propagation formulation and multi-slot time-window constraints, providing optimality certificates as an upper-bound reference. The Tabu Search is augmented with a successive filtering mechanism based on weakness and fitness scores to keep iterations tractable, and with strategic oscillation to systematically escape plateaus by traversing infeasible corridors between feasible solution clusters. Finally, the GRASP metaheuristic is implemented with a scoring function that penalises visit duration in addition to travel time, and a local search phase using prioritised Replace, Insert, and Swap operators with first-improvement Swap to avoid worst-case blow-up. The proposed methods were evaluated using both locally generated datasets and the benchmark dataset of Montemanni and Gambardella commonly used in the OPTW literature.